\SourceLeaf{527}wovon eine leise Andeutung in der Schlussanmerkung der Disquis. Arithm. S. 668 [\textsc{Gauss} Werke B. I. S. 466] gegeben ist.
\par Es ist dies nemlich ein ausreichendes Criterium für den Fall, wo $p$ von der Form $8n+5$ ist.
\par Es sei die Anzahl der Classen, welche die binären Formen in jeder der beiden Gattungen für den Determinant $-p$ bilden $=k$. Der Anfang einer von mir bis zu dem Determinant $-3000$ construirten Tafel steht Disquis. Arithm. p. 520. [art. 303.] Auch ist noch zu bemerken, dass für ein $p$ von der angenommenen Form, allemahl $k=2m+1$ wird, wenn $m$ die Anzahl der Zerlegungen von $p$ in drei positive Quadrate bedeutet (ich sage positiver, um $0$ auszuschliessen), wie \textsc{Legendre} durch Induction gefunden, und in den Disquis. Arithm. zuerst aus der Theorie der \textit{ternären} Formen bewiesen ist. Man hat z. B.
\begingroup\fontsize{8.7}{11}\selectfont\setlength{\arraycolsep}{2pt}
\[
\begin{array}{r*{10}{c}l}
\text{für }p=&5,&13,&29,&37,&53,&61,&101,&109,&149,&157&\text{u. s. w.}\\
k=&1,&1,&3,&1,&3,&3,&7,&3,&7,&3&\\
m=&0,&0,&1,&0,&1,&1,&3,&1,&3,&1&\\[5pt]
&&&\overline{\begin{array}{c}4\\9\\16\end{array}}&\overline{\vphantom{\begin{array}{c}0\\0\\0\end{array}}\hspace{1em}}&\overline{\begin{array}{c}1\\16\\36\end{array}}&\overline{\begin{array}{c}9\\16\\36\end{array}}&\SourceVectorBrace{\begin{array}{ccc}1&4&16\\36&16&36\\64&81&49\end{array}}&\overline{\begin{array}{c}9\\36\\64\end{array}}&\SourceVectorBrace{\begin{array}{ccc}1&4&36\\4&64&64\\144&81&49\end{array}}&\overline{\begin{array}{c}4\\9\\144\end{array}}&
\end{array}
\]
\endgroup
Dies vorausgesetzt, ist allemahl derjenige Werth von $b$, welcher $\equiv\frac12rr\, (\mathrm{mod}.p)$ ist,
\[
{}\equiv2k+a-1\equiv4m+a+1\ (\mathrm{mod}.8)
\]
wodurch das Zeichen von $b$ vollkommen bestimmt ist. Sehen Sie hier 22 Beispiele, indem ich die Ausdehnung der am Schluss der Abhandlung gegebenen Tafel verdopple.
\begin{center}
\begin{tabular}{r|r|r|r||r|r|r|r}
$p$&$k$&$a$&$b$&$p$&$k$&$a$&$b$\\\hline
$5$&$1$&$+1$&$+2$&$181$&$5$&$+9$&$+10$\\
$13$&$1$&$-3$&$-2$&$197$&$5$&$+1$&$-14$\\
$29$&$3$&$+5$&$+2$&$229$&$5$&$-15$&$+2$\\
$37$&$1$&$+1$&$-6$&$269$&$11$&$+13$&$+10$\\
$53$&$3$&$-7$&$-2$&$277$&$3$&$+9$&$+14$\\
$61$&$3$&$+5$&$-6$&$293$&$9$&$+17$&$+2$\\
$101$&$7$&$+1$&$-10$&$317$&$5$&$-11$&$+14$\\
$109$&$3$&$-3$&$+10$&$349$&$7$&$+5$&$+18$\\
$149$&$7$&$-7$&$-10$&$373$&$5$&$-7$&$+18$\\
$157$&$3$&$-11$&$-6$&$389$&$11$&$+17$&$-10$\\
$173$&$7$&$+13$&$+2$&$397$&$3$&$-19$&$-6$\\
\end{tabular}
\end{center}
